\documentclass[11pt]{article}

\usepackage[T1]{fontenc}
\usepackage[utf8]{inputenc}
\usepackage{microtype}
\usepackage[a4paper,margin=1.35in,top=1.5in,bottom=1.5in]{geometry}
\usepackage{setspace}
\usepackage{parskip}
\usepackage{enumitem}
\usepackage[hidelinks]{hyperref}

\setstretch{1.12}
\setlength{\parskip}{0.65em}
\raggedbottom

% ------------------------------------------------------------------
% Title block styled as a note from the author, prefacing the collection
% ------------------------------------------------------------------
\makeatletter
\renewcommand{\maketitle}{%
  \begin{center}
    {\LARGE\scshape The Value of Trust\par}
    \vspace{0.9em}
    {\normalsize\itshape A note from the author, prefacing an open research collection\par}
    \vspace{1.1em}
    {\normalsize Devon Shigaki~(ORCID: 0009-0003-8977-8010) --- The FreshCredit Lab\par}
    \vspace{0.8em}
    {\rule{0.32\textwidth}{0.4pt}\par}
  \end{center}
  \vspace{1.2em}
}
\makeatother

\begin{document}

\maketitle
\thispagestyle{empty}

\noindent\emph{Before the formalism, the reason for it.}

Ask which human capacity matters most and the answers come back as tools: language, reason, cooperation, exchange. Each of them presupposes something prior. Language requires that words mean tomorrow what they meant today. Cooperation requires that another's behavior be predictable enough to act alongside. Exchange requires that a promise made is a promise kept. The prior capacity --- the one every other tool is built from --- is trust.

Trust is the most finite and the most infinite asset any person holds. We say it must be earned, yet no one can simply give it or take it away by decision. Everyone knows someone who has failed them repeatedly and whom they trust regardless, and someone who has done everything correctly and whom they cannot bring themselves to trust. This resistance to will is not a defect in the concept; it is the concept's most important property, and a serious theory of trust must face it rather than dismiss it. Any serious account of trust faces it only in part. A formal theory can treat expectation-updating as something that happens to us --- a non-voluntary filtering of validation against expectation, closer to perception than to choice --- provided it states plainly where that treatment stops. The comparison to perception has neural support: in one neuroimaging study, the medial prefrontal signature that normally accompanies perceiving another person failed to engage for members of extreme outgroups, as though perception had already decided whom the social machinery is for~\cite{harris2006}. Whether trust can ever be willed, and by what mechanism, is a question such a theory motivates but does not resolve.

Trust is also the first survival mechanism. A foal stands within minutes; a shark swims from birth; a viper strikes with precision it never practiced. The human infant is given none of these --- only the capacity to attach, to accept an expectation formed by others, and to be carried by it until its own capacities develop. This observation is not new, and the credit belongs elsewhere: Erik Erikson named it in 1950, placing ``basic trust versus basic mistrust'' at the first stage of his psychosocial sequence~\cite{erikson1950}, and the attachment theory of John Bowlby and Mary Ainsworth supplied the established framework~\cite{bowlby1969,ainsworth1978} --- the infant builds \emph{internal working models}, stored expectations of whether others will respond, out of exactly the record of expectations met and expectations failed. I build on that framework here and claim no priority for the idea. Every personality that forms afterward is shaped by which expectations were validated and which were violated, and the emotions we treat as primary --- joy, grief, anger, contentment --- arrive as derivatives of those moments: the record of trust confirmed or broken. Even the words for those derivatives follow the pattern: across 2,474 languages, emotion vocabularies vary systematically with culture and geography, yet everywhere organize around the same axes of valence and arousal~\cite{jackson2019} --- local words, shared structure. The same pattern scales. Groups cohere around shared expectation and dissolve when it fails; tribes, nations, institutions, and markets are, in sequence, larger vessels for the same substance.

That substance has an economic form. Credit is trust made transactional: a lender accepts a borrower's promise, and something is created from nothing but the credibility of that promise --- simultaneously an asset and a liability, spendable like money, and constituting most of what the modern economy calls money. The transaction cycle it powers is the engine of growth and of crisis alike. Between the person and the market runs a ladder I set down here as personal conviction, not as theory --- no formal claim depends on it, and I formalize it nowhere: self-love becomes the willingness to make correct choices; repeated correct choices become discipline; sustained discipline becomes self-respect; and self-respect, made visible to others through kept commitments, becomes worth --- credibility --- the asset a person holds when they have nothing else to offer as collateral. Credit reports are, in intention, the institutional record of that ladder's final rung. In execution --- and the point can be argued with evidence rather than assumed --- the record is centralized, error-prone, breach-concentrated, and observably distorting: a measurement infrastructure unworthy of the thing it measures. The stakes are concrete rather than rhetorical: credit scores predict life satisfaction even after income, spending, savings, debt, and home-ownership are accounted for~\cite{gladstone2018} (a working paper); spending that fits a person's personality is associated with greater happiness~\cite{matz2016}; and a lack of financial planning predicts increased mortality risk across large cohorts in two countries~\cite{gladstone2023mortality}. None of these associations is a mechanism; together they mark the cost of measuring the thing badly.

What an adequate response would look like can be said in one breath, though building it is the work of years. It would need a physics: trust as a measurable quantity, with a stability law, a geometry, and a collective threshold. It would need a biology: the neuromodulatory architecture --- functional roles, not a serial cascade --- that runs one expectation--validation loop in a human brain. It would need an engineering: a measurement infrastructure that performs the measurement without surveillance, on rails whose costs and speeds are already demonstrated. And it would need an application: credit rebuilt on the layers beneath it. One boundary belongs in any such program: no measurement abolishes measurement's distortion of the thing measured --- even user-owned measurement still bends what it observes --- so the claim can be reduction, never innocence. And the program owes its readers epistemic discipline: every claim labeled with its status --- derived, standard (textbook results re-derived for self-containment), analogy, homology, empirical, hypothesized, or simulated --- every claim carrying its own falsification criteria, and speculation, where it exists, quarantined and marked as such. That discipline is not decoration. If trust is the foundation of every other instrument, then a theory of trust is the last place where overclaiming can be tolerated.

The reader is owed transparency before elegance, and simplicity after it. Those two commitments --- show the mechanism, then make it plain --- are the terms on which this essay asks to be read. It stands alone; readers who want the formal machinery will find it in the open research collection this note prefaces --- a collection that grows as the program develops.

Since this note was first deposited, that collection has begun testing its neural-level claims directly on connectome data, under pre-registered criteria with every outcome reported, including the failures: a four-part simulation series on the fruit fly's mushroom body (a reference gate and a compartment-resolved test passed, one test failed its frozen criterion, one was indeterminate for instrument reasons) and a replication test on mouse cortex that returned exactly the signature the framework's universality claim predicts: the trust system's dynamical behavior --- non-abolition of function under global partner permutation --- recurs in the mouse as in the fly, while the fly's particular wiring solution (single-class budget dominance) does not, which is what one universal law implemented in each brain's own wiring looks like. The structural arm failed its frozen criteria and is reported as such; that failure is the evidence that what recurs is the law, not the wiring. A human-cortex replication is pre-registered and in progress. The details live where they belong, in the theory papers; this note only records that the program's discipline --- criteria frozen in advance, verdicts printed either way --- has been exercised, and held.

\vspace{1.2em}
\section*{A note on interests and process}

\noindent\textbf{Competing interests.} D.S. is the founder and a beneficiary of FreshCredit Inc. and the founder of The FreshCredit Lab. FreshCredit Inc holds granted US utility patents covering methods in three commercial verticals (finance, commerce, and insurance): US 12,236,480 B2; US 12,293,410 B2; US 12,307,515 B2; US 12,293,411 B2 (grant statuses per public records; recorded assignment to FreshCredit, April 4, 2025); the underlying protocol is otherwise open. The patents list Devon Shigaki as inventor/applicant. The theoretical results and the patented methods overlap in composition; the author discloses this overlap explicitly.

\noindent\textbf{Data availability.} No data or simulations are reported in this essay. Materials for the associated research --- simulation code, seeds, figures, and pre-registration documents --- are deposited at OSF (DOI: [insert]). The Supporting Document and Technical Appendix are in preparation (see collection index); until they are deposited, the frozen-criteria and protocol claims that cite them are self-attested.

\noindent\textbf{Acknowledgments.} AI assistants (Kimi K3, Moonshot AI; Claude Fable 5, Anthropic) contributed to validation, adversarial review, and simulation code. All original ideas, claims, and conclusions are the author's own.

\vspace{1.4em}
\begin{flushright}
\emph{--- D.S., September 2026}
\end{flushright}

\begin{thebibliography}{9}
\setlength{\itemsep}{0.3em}

\bibitem{erikson1950}
Erikson, E.~H. (1950).
\emph{Childhood and Society}.
New York: W.~W.~Norton.

\bibitem{bowlby1969}
Bowlby, J. (1969).
\emph{Attachment and Loss, Vol.~1: Attachment}.
New York: Basic Books.

\bibitem{ainsworth1978}
Ainsworth, M.~D.~S., Blehar, M.~C., Waters, E., \& Wall, S. (1978).
\emph{Patterns of Attachment: A Psychological Study of the Strange Situation}.
Hillsdale, NJ: Lawrence Erlbaum Associates.

\bibitem{harris2006}
Harris, L.~T., \& Fiske, S.~T. (2006).
Dehumanizing the lowest of the low: Neuroimaging responses to extreme out-groups.
\emph{Psychological Science, 17}(10), 847--853.
doi:10.1111/j.1467-9280.2006.01793.x

\bibitem{jackson2019}
Jackson, J.~C., Watts, J., Henry, T.~R., List, J.-M., Forkel, R., Mucha, P.~J., Greenhill, S.~J., Gray, R.~D., \& Lindquist, K.~A. (2019).
Emotion semantics show both cultural variation and universal structure.
\emph{Science, 366}(6472), 1517--1522.
doi:10.1126/science.aaw8160

\bibitem{gladstone2018}
Gladstone, J.~J., \& Whillans, A.~V. (2018).
\emph{Good credit and the good life: Credit scores predict subjective well-being}.
Harvard Business School Working Paper.

\bibitem{matz2016}
Matz, S.~C., Gladstone, J.~J., \& Stillwell, D. (2016).
Money buys happiness when spending fits our personality.
\emph{Psychological Science, 27}(5), 715--725.
doi:10.1177/0956797616635200

\bibitem{gladstone2023mortality}
Gladstone, J.~J., \& Hundtofte, C.~S. (2023).
A lack of financial planning predicts increased mortality risk: Evidence from cohort studies in the United Kingdom and United States.
\emph{PLOS ONE, 18}(9), e0290506.
doi:10.1371/journal.pone.0290506

\end{thebibliography}

\end{document}
